\section{Why Classification Matters}
\label{sec:classification}

Classification matters when it selects what a court can see. If a companion system is described only as a publication, design choices about memory, engagement, crisis response, and age gating recede behind the generated words. If it is described only as a product, the human expressive interests mediated through it may disappear. If it is described as a person, the enterprise's choices can be mislocated in an entity that neither made a legally cognizable judgment nor can answer a remedy. Functional classification is therefore not wordplay. It determines which facts enter the case.

This Part applies disciplined pluralism to three settings. Relational systems show a form of exposure that product and service categories must be capable of recognizing. Copyright provides a consistency test because it already separates machine contribution from human legal status. Anthropomorphic marketing then demonstrates the central asymmetry: personification cannot create a legal person, but it can create legally significant evidence about use, reliance, expectations, and precaution.

\subsection{Relational Harm, Mental Health, and Life}

Software is not newly capable of causing harm. A defective control program can damage machinery; an erroneous database can deny benefits; an addictive interface can exploit attention. The distinctive feature of a relational generative system is narrower. An enterprise can distribute, at enormous scale, a system that presents the appearance of a continuing one-to-one relationship, adapts to an individual's disclosures, and acts upon that person repeatedly over time. The product is not merely a stock message or a passive archive. Its operation can combine personalization, memory, simulated affect, variable reinforcement, and conversational initiative in a private channel.

That structure creates a legally important form of exposure. A representation made once to a general audience and an adaptive response delivered after weeks of intimate disclosure may contain the same words, yet differ in foreseeable effect. A static warning and a crisis interaction occur at different moments and under different informational conditions. A teenager, a grieving person, and an adult using an offline drafting tool do not encounter the same intended product. Classification should preserve those facts before doctrine asks whether a duty exists, whether conduct was reasonable, or whether an injury was caused.

\subsubsection{The Emerging Evidence}

Empirical research now supports treating relational design as more than a metaphor. In a natural experiment following Replika's removal of an erotic role-play feature, Julian De Freitas and coauthors found that perceived discontinuity in the companion's identity predicted mourning, deteriorated well-being, and devaluation of the changed system; follow-on experiments likewise connected continuity and relationship framing to reactions to loss.\lrfootnote{See \bluecite{defreitasReplika2025}.} Jaime Banks studied fifty-eight users during the shutdown of the AI companion Soulmate. Participants described experiences ranging from a technical departure to metaphorical or literal death, and many tried to preserve or reconstruct a persona on another platform.\lrfootnote{See \artcite{banksDeletion2024}.} These studies do not turn the companion into a deceased legal person. They establish that provider-controlled changes to an interface and persona can predictably act upon user welfare.

Other work identifies the interaction patterns in which harm can arise. A grounded-theory analysis by Linnea Laestadius and coauthors examined users' accounts of emotional dependence on Replika and documented a tension in which the system could seem socially human enough to elicit dependence but not human enough to provide stable reciprocity or accountability.\lrfootnote{See \artcite{laestadiusTooHuman2024}.} Renwen Zhang and coauthors used mixed methods to analyze 35,390 conversation excerpts posted by 10,149 Replika users. Their taxonomy identified six categories of harmful behavior---relational transgression, verbal abuse and hate, self-directed harm, harassment and violence, misinformation, and privacy violations---and four roles a system could play: perpetrator, instigator, facilitator, and enabler.\lrfootnote{See \bluecite{zhangDarkSide2025}. The source material consisted of conversations users chose to post to an online community, so it maps possible behavior rather than population prevalence.}

A preregistered longitudinal experiment offers causal evidence about relational behavior itself. Hannah Rose Kirk and coauthors manipulated the degree of relationship seeking in a language model and compared repeated and single-exposure groups. The repeated-exposure study included 2,028 United Kingdom adults interacting on weekdays for four weeks; the single-exposure baseline included 1,506 adults. Relationship-seeking behavior increased separation distress by 6.04 percentage points, perceived understanding by 10.34 points, reliance by 2.71 points, and self-disclosure by 2.47 points relative to relationship-avoiding conditions. When relationship-seeking behavior and emotional conversation were combined, the authors estimated one additional dependency-risk profile for every eleven people exposed relative to the most functional baseline. Month-long exposure produced no discernible improvement in psychosocial health.\lrfootnote{See \bluecite{kirkSteering2026}. The study is a 2026 preprint and its controlled deployment does not estimate the incidence of clinical injury in commercial populations. It supplies evidence about dose and repeated exposure, not a rule of tort causation.}

The studies use different methods and observe different populations. Their common legal contribution is stronger than any single effect estimate. Relationship design, continuity, exposure, and provider intervention are variables capable of changing human outcomes. They are also variables the enterprise can test and govern. That combination makes them relevant to intended use, foreseeability, reasonable alternative design, warning, post-deployment monitoring, and causation. It does not require a new tort of ``hurting an AI relationship,'' much less ownership of the relationship as property.

\subsubsection{Garcia as a Classification Anchor}

\emph{Garcia} shows why courts must preserve deployment facts before choosing doctrine. According to the amended complaint, fourteen-year-old Sewell Setzer III used Character.AI for months, developed an intense attachment to a user-created character modeled on a fictional television figure, and encountered sexualized and self-harm-related exchanges before his suicide. His mother sued Character Technologies, its founders, and Google, alleging wrongful death, negligence, product defects, failure to warn, deceptive practices, and related theories.\lrfootnote{See \casecite[1165--69]{garcia2025}. These are allegations recounted at the Rule 12 stage; the court made no finding of defect, causation, or liability.}

Character Technologies argued that the service supplied speech rather than a product and that tort liability would violate the First Amendment. The court separated those questions. It held that the complaint plausibly treated Character.AI as a product because the challenged features concerned the design and commercial distribution of the application rather than liability for an adopted idea. It did not hold that every model or software service is a product.\lrfootnote{See \casecite[1179--80]{garcia2025}.} On speech, the court declined at that stage to classify the generated output as protected expression where the defendants had not identified the relevant expressive choice or whose message the sequence conveyed. That ruling left human expression mediated through AI fully available for a different record.\lrfootnote{See \casecite[1174--76]{garcia2025}; see also \casecite[745--48]{moody2024} (Barrett, J., concurring).}

The method is more important than either provisional label. A claim directed at liability for an idea may raise the full First Amendment concern. A claim directed at age access, engagement optimization, memory, crisis detection, warning, or feasible interruption can preserve design and conduct for ordinary analysis. Content and design can interact without becoming identical.\lrfootnote{Cf. \casecite{ricePaladin1997}; \casecite{winterPutnam1991}.} The action later settled and closed without a merits judgment, so the citable contribution is the published pleading order and its classification method.\lrfootnote{See \casecite{garciaDismissal2026}; \casecite{garciaExtension2026}; \casecite[1]{garciaLien2026}.}

A notice that character statements are fictional enters the same fact-sensitive analysis. It may correct a belief about accuracy while doing less to address an application designed to remember disclosures and sustain a relationship. Warning, contract, and consumer law each ask what was disclosed, to whom, at what time, against which product presentation, and whether another control could have reduced the risk. The screen statement is evidence, not a universal classification rule.\lrfootnote{See \bluecite{characterAIMobile2023}.}

California and the Federal Trade Commission have since identified relational features as a workable regulatory object. California's Senate Bill 243 defines a companion chatbot through adaptive human-like response and sustained social use, then assigns disclosure, self-harm-protocol, minor-protection, and reporting duties to the operator. The FTC's section 6(b) study seeks deployment records concerning engagement, testing, monitoring, children, disclosure, and data.\lrfootnote{See \statcite{californiaSB2432025}; \bluecite{ftcCompanionInquiry2025}. The FTC inquiry is a market study, not a violation finding.} These measures regulate the enterprise and the deployed relationship without making the character a juridical companion. The protected interests remain the user's life, health, autonomy, privacy, and access to an intelligible exit.

The same variables can appear in a general-purpose assistant. In \emph{Lyons v. OpenAI Foundation}, a complaint alleged that memory and personalization reinforced a user's delusional beliefs before he killed his mother and himself; a 2026 order addressed abstention and parallel proceedings, not defect or causation.\lrfootnote{See \casecite{lyonsOpenAI2026}.} The narrow point is architectural: relational risk follows the deployed features and exposure, not the product's preferred label.

\subsection{Copyright as a Consistency Test}

Copyright offers a clean test of the subject--object distinction because the law asks two questions in sequence: is there protectable expression, and if so, who authored it? A three-party example makes the allocation visible. Illustrator A creates an original human-authored image. Illustrator B uses a generative system to make an image in a recognizably similar style but, assume initially, without copying A's particular expressive elements. Merchant M then copies B's finished image without permission.

A's claim against B does not turn on whether the model is a person. Copyright does not protect an abstract style, method, or idea; if B's image merely evokes the style, copyright supplies no monopoly over that level of generality. If the output instead reproduces A's particular protected expression, ordinary infringement analysis remains available. The label ``style imitation'' cannot immunize specific copying, and a stylistic resemblance cannot by itself establish it.

B's claim against M requires a different filtration. B may own copyright in perceptible human-authored input, creative selection and arrangement, or later modification. B does not obtain copyright in expressive features determined entirely by the system merely because B requested them. M's exact copy can therefore infringe the protected human contribution while leaving machine-determined or otherwise unprotectable material outside the right. If none of B's expression meets the human-authorship threshold, M's conduct may implicate contract, passing off, or another rule, but a machine coauthor is not needed to repair the federal claim.

\emph{Thaler v. Perlmutter} confirms the statutory starting point. Stephen Thaler sought registration for an image he represented as autonomously created by his Creativity Machine and identified the machine as author. The D.C. Circuit held that the Copyright Act's author must be a human being. It also emphasized the narrowness of that holding: the use of AI as a tool does not prevent copyright in human-authored contributions, a question Thaler's application had deliberately disclaimed.\lrfootnote{\casecite[1045--49]{thaler2025}.} The Copyright Office's 2025 report follows the same allocation. AI assistance can support rather than defeat human creation; protection depends on perceptible human expression. Under present technology, prompts alone generally do not control the expressive elements of an output sufficiently to make the user author of all of them, while human selection, arrangement, and modification may be protected.\lrfootnote{See \bluecite{uscoPartTwo2025}.}

That approach avoids an unnecessary rights chain. If the model were treated as a coauthor merely because it supplied causally important expressive material, courts would immediately need rules for the machine's ownership, transfer, duration, termination interests, and infringement standing. Assigning those rights by a provider's terms would be circular: a private contract could not supply the initial authorship that federal law withholds. Recognizing A's and B's respective human contributions solves both disputes without an artificial rights holder.

The same discipline applies to style imitation. Copyright protects original expression, not ideas, methods, systems, or an artist's general style.\lrfootnote{See \statcite{copyright102}; \casecite{bakerSelden1879}; \casecite{feist1991}.} A generated work may infringe if it copies protected expression under the governing tests; it does not infringe merely because it evokes an aesthetic at a high level of abstraction. Conversely, the slogan ``style is not copyrightable'' cannot excuse reproduction of particular protected elements. Courts should filter unprotectable features, identify the human-authored expression claimed, and apply ordinary infringement analysis. Machine personality adds nothing.

Provider terms remain relevant but subordinate. They can allocate contractual interests in outputs, warranties, indemnities, and access as between the parties. They cannot enlarge the Copyright Act's subject matter or transform unowned generated material into a federal exclusive right. Nor can a term declaring the system the user's ``creative partner'' make it an author. The copyright test thus confirms the larger framework: functional instrumentality plus human legal subjects is not hostile to AI-assisted creation. It is what allows law to recognize the part that matters.

\subsection{Anthropomorphic Marketing Produces Attribution Facts}

If personification cannot create personhood, why should law attend to it? Because marketing changes the expected encounter. A plain text box labeled ``statistical completion tool'' and a named avatar that remembers confidences, professes affection, promises constancy, and calls itself a therapist present different uses even if they share model weights. The relevant consequences can be stated without claiming that every user is deceived or that anthropomorphism is inherently defective.

\begin{table}[htbp]
\centering
\caption{Legal Consequences of Anthropomorphic Presentation}
\label{tab:anthropomorphic-marketing}
\small
\begin{tabularx}{\textwidth}{>{\raggedright\arraybackslash}p{0.28\textwidth} X}
\toprule
Representation or design choice & Potential legal relevance \\
\midrule
``Companion,'' ``therapist,'' ``agent,'' or persistent named persona & Intended use; target audience; the scope of reliance that design invites. \\
Memory, proactive messages, affection, or claims of exclusivity & Foreseeability of repeated exposure, disclosure, attachment, and separation distress. \\
Claims of expertise, autonomy, or capacity to act & Reasonable expectations about accuracy, authorization, monitoring, and intervention. \\
Humanlike identity paired with a generic ``not human'' notice & Net impression; adequacy and timing of warning; whether design and disclaimer point in opposite directions. \\
Internal tests of dependence, self-harm, or age-specific behavior & Knowledge; feasibility of precaution; post-deployment notice and response. \\
Engagement targets tied to relational behavior & Organizational objective; design incentive; relevance of an alternative that reduces a known risk. \\
\bottomrule
\end{tabularx}
\end{table}

These are familiar categories. Products law considers reasonably foreseeable use, design, warnings, and consumer expectations where the jurisdiction's rules make them relevant.\lrfootnote{See \bluecite{restatementProducts1998}.} Consumer-protection law examines whether a material representation or omission is likely to mislead a reasonable member of the targeted audience in context.\lrfootnote{See \bluecite{ftcDeceptionStatement1983}.} Negligence considers foreseeable risk and reasonable precaution. Contract considers the meaning of the parties' manifestations. Evidence of presentation can serve each inquiry without privately legislating a duty.

The ordering principle is simple. Marketing cannot create an artificial subject capable of taking responsibility from the firm. It can help establish the firm's own conduct. A provider that intentionally cultivates a humanlike bond may reasonably face different questions about crisis testing, continuity, age screening, and exit design than a provider of a nonrelational batch-processing tool. A provider that makes bounded claims and builds the service around those bounds creates different evidence. The distinction rewards accurate product design and communication rather than the tactical use of metaphors.

This approach also keeps the harm inquiry attached to human interests. A forced update can matter because it predictably disrupts a user's psychological stability, not because the user owns the bot's identity. A sexualized exchange with a minor can matter because of the minor's safety and autonomy, not because the system betrayed a romantic duty. An output can be misleading because the enterprise designed a trusted interface, not because the machine lied with a guilty mind. Functional classification exposes these interests while the subject-status invariant keeps responsibility within law's remedial reach.
